\RequirePackage[T1]{fontenc}
\documentclass[10pt,journal]{IEEEtran}
\usepackage[letterpaper,top=1in,bottom=1in,left=0.7in,right=0.7in,columnsep=0.2in]{geometry}

\usepackage{iftex}
\ifPDFTeX
\usepackage[utf8]{inputenc}
\usepackage[T1]{fontenc}
\fi
\usepackage{amsthm}
\usepackage{amsmath, bm}
\usepackage{newtxtext,newtxmath}
\usepackage{textcomp}
\interdisplaylinepenalty=2500
\usepackage{pgf}
\usepackage{pgfplots}
\usepackage{tikz}
\usetikzlibrary{arrows.meta,fit,backgrounds,positioning,calc}
\usepackage{xcolor}
\usepackage{float}
\usepackage{mathrsfs}
\usepackage{graphicx}
\usepackage{svg}
\usepackage{placeins}
\usepackage{array, booktabs, makecell, adjustbox}
\usepackage{mathtools}
\usepackage{nccmath, relsize}
\usepackage{enumitem}
\usepackage{flushend}
\usepackage{flafter}
\usepackage{siunitx}
\usepackage[final]{microtype}
\usepackage{etoolbox}

\DeclareRobustCommand{\vct}[1]{\bm{#1}}
\DeclareRobustCommand{\mat}[1]{\bm{#1}}
\DeclareRobustCommand{\matg}[1]{\bm{#1}}
\newcommand{\Quat}{\mathbb{H}}
\newcommand{\Sph}{\mathbb{S}}
\newcommand{\SO}{\mathrm{SO}}
\newcommand{\SE}{\mathrm{SE}}
\newcommand{\quat}[1]{\vct{#1}}
\newcommand{\Rot}{\bm{\mathcal{R}}}
\newcommand{\qconj}[1]{\overline{#1}}
\newcommand{\qreal}[1]{\operatorname{Re}\! \left(#1\right)}
\newcommand{\qimag}[1]{\operatorname{Im}\! \left(#1\right)}

\usepackage{cite}
\usepackage[breaklinks=true,hidelinks]{hyperref}
\usepackage{bookmark}
\usepackage[nameinlink,capitalise,noabbrev]{cleveref}

\sisetup{detect-weight=true,detect-family=true,per-mode=symbol}
\setlist[description]{font=\footnotesize\bfseries,labelsep=0.6em,leftmargin=1.8em}

\pdfstringdefDisableCommands{%
    \def\SI#1#2{#1 #2}%
    \def\bm#1{#1}%
    \def\Skew#1{[#1]_\text{x}}%
    \def\Id{I}%
    \def\pct#1{#1\%}%
    \def\circ{deg}%
    \def\degree{deg}%
    \def\propto{∝}%
    \def\sigma{sigma}%
    \def\tau{tau}%
    \def\Phi{Phi}%
    \def\omega{omega}%
    \def\_{}%
    \def\cdot{}%
    \def\;{ }%
    \def\I{I}%
    \def\quat#1{#1}%
    \def\SO{SO}%
    \def\SE{SE}%
    \def\Quat{H}%
    \def\Sph{S}%
}

\AtBeginEnvironment{equation}{\small}
\AtBeginEnvironment{align}{\small}
\providecommand{\mathdefault}[1]{#1}

\setcounter{topnumber}{5}
\setcounter{dbltopnumber}{5}
\renewcommand{\topfraction}{0.95}
\renewcommand{\dbltopfraction}{0.95}
\renewcommand{\textfraction}{0.05}
\renewcommand{\floatpagefraction}{0.9}
\renewcommand{\dblfloatpagefraction}{0.9}
\newcommand{\IncludePGF}[1]{\resizebox{\linewidth}{!}{\input{#1}}}

\ifPDFTeX
\DeclareUnicodeCharacter{2013}{\textendash}
\DeclareUnicodeCharacter{2014}{\textemdash}
\DeclareUnicodeCharacter{2018}{\textquoteleft}
\DeclareUnicodeCharacter{2019}{\textquoteright}
\DeclareUnicodeCharacter{201C}{\textquotedblleft}
\DeclareUnicodeCharacter{201D}{\textquotedblright}
\DeclareUnicodeCharacter{2212}{\textminus}
\fi

\pdfstringdefDisableCommands{%
    \def\vct#1{#1}%
    \def\mat#1{#1}%
    \def\matg#1{#1}%
    \def\I{I}%
}

\newtheoremstyle{bodysmall}{3pt}{3pt}{\normalfont\small}{0pt}{\bfseries\small}{.}{0.5em}{}
\theoremstyle{bodysmall}
\newtheorem{lemma}{Lemma}
\newtheorem{theorem}{Theorem}
\newtheorem*{theorem*}{Theorem}
\theoremstyle{remark}
\newtheorem{remark}{Remark}

\newcommand{\Skew}[1]{\left[#1\right]_\times}
\newcommand{\Id}{\mat{I}}
\newcommand{\pct}[1]{#1\%}
\newcommand{\meanstd}[2]{#1 \pm #2}
\newcommand{\ci}[2]{#1\,(\,#2\,\text{CI}\,)}
\newcommand{\Ito}{\mathrm{It\hat{o}}}
\newcommand{\ItoInt}[2]{\int_0^{#1} #2\, dW_s^{\Ito}}
\DeclareMathOperator{\proj}{\mathcal{P}}
\DeclareMathOperator{\clip}{clip}
\DeclareMathOperator{\wrap}{wrap}
\DeclareMathOperator{\sgn}{sgn}
\DeclareMathOperator{\blkdiag}{blkdiag}
\DeclareMathOperator{\diag}{diag}
\DeclareMathOperator{\scomp}{scomp}
\newcommand{\set}[1]{\mathcal{#1}}
\newcommand{\R}{\mathbb{R}}
\newcommand{\C}{\mathbb{C}}
\newcommand{\N}{\mathbb{N}}
\newcommand{\Z}{\mathbb{Z}}
\newcommand{\E}{\mathbb{E}}
\DeclareMathOperator{\Var}{Var}
\DeclareMathOperator{\Cov}{Cov}
\newcommand{\ts}[1]{\text{#1}}
\newcommand{\fitcol}[1]{\adjustbox{max width=\columnwidth}{\ensuremath{\displaystyle #1}}}

\makeatletter
\@ifpackageloaded{cite}{\renewcommand{\citepunct}{,\penalty\citepunctpenalty\,}\renewcommand{\citedash}{--}}{}
\makeatother
\emergencystretch=2em
\allowdisplaybreaks

\title{Adaptive OU-Regularized Quaternion MEKF for Marine Wave-Motion Estimation}
\author{Mikhail S. Grushinskiy\thanks{M. S. Grushinskiy is an independent
researcher with the Bareboat Necessities project, Fair Lawn, NJ, USA
(e-mail: mgrouch@users.github.com).}}
\pgfplotsset{compat=1.18}

\input{w3d-editorial-tables.tex-part}

\begin{document}

% Result macros share the authoritative generated evidence with the tables.
\IfFileExists{w3d-ou-validation-macros-generated.tex-part}{%
  \input{w3d-ou-validation-macros-generated.tex-part}%
}{}
\IfFileExists{w3d-ou-robustness-macros-generated.tex-part}{%
  \input{w3d-ou-robustness-macros-generated.tex-part}%
}{}
\providecommand{\OURobustnessNominalMean}{--}
\providecommand{\OURobustnessSigmaSpan}{--}
\providecommand{\OURobustnessRSSpan}{--}
\providecommand{\OURobustnessCoupledSigmaSpan}{--}
\providecommand{\OURobustnessCoupledTauSpan}{--}
\providecommand{\OUValidationBootstrapMethod}{a nonparametric percentile bootstrap of the paired mean}
\providecommand{\OUValidationBootstrapResamples}{10{,}000}
\providecommand{\OUValidationBootstrapRNG}{NumPy PCG64}
\providecommand{\OUValidationStatsSeed}{--}
\providecommand{\OUValidationPairedComparisons}{--}
\providecommand{\OUValidationThreeDScenarios}{--}
\providecommand{\OUValidationThreeDHigherCount}{--}
\providecommand{\OUValidationThreeDUnresolvedCount}{--}
\providecommand{\OUValidationNormalizedTLow}{--}
\providecommand{\OUValidationNormalizedTHigh}{--}
\providecommand{\OUValidationNormalizedTStatistic}{--}
\providecommand{\OUValidationNormalizedTP}{--}
\providecommand{\OUValidationNormalizedSignNegative}{--}
\providecommand{\OUValidationNormalizedSignP}{--}
\providecommand{\OUValidationNormalizedRandomizationP}{--}
\providecommand{\OUValidationNormalizedRandomizationPatterns}{--}

\maketitle

\begin{abstract}
Low-cost inertial wave-motion reconstruction must suppress drift without
attenuating the desired vessel displacement. We develop OU--III, a
21-state quaternion multiplicative extended Kalman filter coupling attitude
and sensor biases to an Ornstein--Uhlenbeck-driven velocity, displacement,
integral-displacement, and acceleration chain. A zero pseudo-measurement on
the integral state regularizes low frequencies. Measurement-only estimates
of wave period and band-limited acceleration schedule the OU prior and a
SpectralMSE regularizer derived from a physical drift--distortion balance,
including pseudo-update cadence. In a paired ten-seed study, stationary
JONSWAP vertical RMS is
$\OUValidationOUIIINormalizedMean\pm\OUValidationOUIIINormalizedStd\%$ of
incident significant wave height for OU--III versus
$\OUValidationOUIINormalizedMean\pm\OUValidationOUIINormalizedStd\%$ for
OU--II. The paired difference is $\OUValidationNormalizedDifference$
percentage points, with bootstrap 95\% interval
$[\OUValidationNormalizedDifferenceLow,\OUValidationNormalizedDifferenceHigh]$.
A four-arm ablation attributes most off-nominal improvement to adaptive
integral regularization rather than OU-prior adaptation alone. Comparisons,
changing-sea tests, and vibration failure/mitigation studies characterize
tradeoffs. Evidence is limited to synthetic fixed-heading motion from one
estimated vessel response; displacement is relative wave motion, not
absolute navigation position.
\end{abstract}

\begin{IEEEkeywords}
Attitude estimation, heave, inertial sensors, Kalman filtering, marine motion,
Ornstein--Uhlenbeck process.
\end{IEEEkeywords}

\input{w3d-intro.tex-part}
\input{w3d-related-work-modern.tex-part}
\input{w3d-nomenclature.tex-part}
\input{w3d-architecture.tex-part}
\input{w3d-state.tex-part}
\input{w3d-meas.tex-part}
\input{w3d-proc-cont.tex-part}
\input{w3d-lti-discrete.tex-part}
\input{w3d-analytic-coeff.tex-part}
\input{w3d-kalm-up.tex-part}

\input{w3d-adaptation-motivation.tex-part}
\input{w3d-adaptation-cadence-interpretation.tex-part}
\input{w3d-reduced-mse-envelope.tex-part}
\input{w3d-reduced-mse-validation.tex-part}
\input{w3d-adaptation-deployed-law.tex-part}
\input{w3d-rs-anisotropy-design.tex-part}
\input{w3d-fus-methods.tex-part}
\input{w3d-adaptation-literature-context.tex-part}

\input{w3d-initialization-overview.tex-part}
\input{w3d-mag-hard-iron.tex-part}

\input{w3d-sim-charts.tex-part}
\input{w3d-results.tex-part}

\input{w3d-post-results-investigations.tex-part}
\input{w3d-conclusion-summary.tex-part}

\FloatBarrier
\appendices
\section{Detailed Numerical Results}
\label{sec:detailed-numerical-results}
The following tables retain the exact numerical results accompanying the
main-text charts. They use the original scoring windows and aggregation;
no numerical entries are removed or recomputed.
\OUPrintAppendixTables
\FloatBarrier

\section*{Acknowledgments}
ChatGPT and Claude assisted with code generation, mathematical drafting, and editorial preparation under author supervision. The author is responsible for the method, validation, and final manuscript.

\section*{Disclosure Statement}
The author declares no conflicts of interest and received no funding for this research.

\nocite{Kalman1960_Filtering,Jazwinski1970_Filtering,Maybeck1979_StochasticModels,Markley2003AttitudeError,UhlenbeckOrnstein1930_OU}
\bibliographystyle{IEEEtran}
\bibliography{w3d,w3d-modern,w3d-iss,w3d-baselines}
\end{document}
